\section{Empirical Results: Breaking the Exploration Bottleneck}
\label{ch:experiments}

This chapter presents the experimental setup, the environments used for evaluation, and the results obtained. We specifically focus on validating the hypothesis that a Meta-Controller LLM can effectively resolve exploration bottlenecks in a sparse Multi-Agent Reinforcement Learning (MARL) environment through dynamic reward shaping, and subsequently, how Hierarchical Reinforcement Learning (HRL) further improves these capabilities.

\subsection{Experimental Setup}
\label{sec:exp_setup}

To rigorously evaluate the proposed architectures, we compare distinct configurations within our multi-agent resource-gathering environment, structured into two parts. 

For the \textbf{Dynamic Reward Shaping} experiments, we evaluate:
\begin{enumerate}
    \item \textbf{Sparse Baseline:} A standard MAPPO implementation~\parencite{Yu_Velu_Vinitsky_Gao_Wang_Bayen_Wu_2022} relying purely on sparse environmental rewards (+1 for completing milestones).
    \item \textbf{LLM Dynamic (Gemini 2.5 Flash Lite):} The proposed ``Mixing Board'' architecture, utilizing Google's Gemini 2.5 Flash Lite model via API as the Commander to dynamically shift potential-based reward weights. 
    \item \textbf{LLM Dynamic (Local Qwen 2.5 7B):} An equivalent dynamic architecture utilizing a localized, smaller parameter model (Qwen 2.5 7B) to observe the efficacy of smaller, self-hosted LLMs. 
\end{enumerate}

For the \textbf{Hierarchical Reinforcement Learning (HRL)} experiments, we evaluate three configurations:
\begin{enumerate}
    \item \textbf{No-LoRA Base (Qwen 2.5 7B):} The un-finetuned base model serving as the HRL Meta-Controller (the Commander).
    \item \textbf{QLoRA with Reflection (Qwen 2.5 7B):} A QLoRA fine-tuned model (as detailed in Section~\ref{sec:qlora_pipeline}, the Qwen 2.5 7B model was fine-tuned using a 4-bit NF4 quantization pipeline on 1,200 synthetic oracle examples) with counterfactual reflection triggered on option staleness.
    \item \textbf{QLoRA without Reflection (Qwen 2.5 7B):} The same fine-tuned model, with the reflection mechanism disabled.
\end{enumerate}

All runs were executed with identical MAPPO hyperparameters (clipping $\epsilon=0.2$, GAE\footnote{Generalized Advantage Estimation} $\lambda=0.95$, $\gamma=0.99$, across 128 parallel environments).

% ==========================================
% PART 1: DYNAMIC REWARD SHAPING
% ==========================================
\subsection{Part 1: Dynamic Reward Shaping}
\label{sec:exp_dynamic}

The most crucial metric validating the LLM Dynamic architecture is the progression through the environment's strict dependency Directed Acyclic Graph (DAG). 

\begin{figure}[H]
    \centering
    \includegraphics[width=\textwidth]{../../plots/baseline_vs_llmdynamic_results/subtasks_comparison.png}
    \caption{Subtask Completion Rates: Comparison between the Sparse Baseline and the LLM Dynamic models over 2000 training updates.}
    \label{fig:subtasks_comparison_dynamic}
\end{figure}

As shown in \cref{fig:subtasks_comparison_dynamic}, all three configurations conquer the early-game tasks (Wood, Stone) almost immediately. The Pickaxe subtask follows shortly after, with the Baseline reaching it at step~1.1M and Gemini at step~720K.

However, the Baseline completely collapses at the \emph{Iron} bottleneck. Without any dynamic reward shaping to guide exploration, the Baseline workers never discover how to consistently mine Iron (0\% final rate), resulting in a total failure to learn all downstream tasks (Sword, Armor, Bridge). 

In contrast, both LLM-guided runs break through the Iron bottleneck, but at remarkably different speeds. \textbf{Qwen} unlocks Iron at step~10M, subsequently discovering Sword and Armor at essentially the same step (10.1M). This rapid cascade demonstrates the snowball effect: once a critical bottleneck is broken, downstream tasks are unlocked quickly. \textbf{Gemini 2.5 Flash Lite} unlocks Iron much later, at step~40.7M (approximately 4$\times$ slower than Qwen). However, once unlocked, the same cascade occurs: Sword and Armor appear at step~41.3M.

Neither run achieved Enemy defeats (0\% for all three configurations), indicating that the combat subtask represents a qualitatively different challenge beyond reward gravity, likely requiring specialised action coordination that the current shaping signal alone cannot provide.

\subsubsection{LLM Interventions in Action}
\label{sec:exp_interventions}

To understand exactly \emph{how} the LLMs broke the bottleneck, we analyse the evolution of the reward weights. The weights shown in \cref{fig:llm_weights_evolution} are the \emph{post-softmax} weights actually applied to the potential-based shaping.

\begin{figure}[H]
    \centering
    \begin{subfigure}{0.8\textwidth}
        \includegraphics[width=\textwidth]{../../plots/baseline_vs_llmdynamic_results/llm_weights_evolution_gemini.png}
        \caption{Gemini 2.5 Flash Lite Meta-Controller (198 interventions)}
    \end{subfigure}
    \vskip\baselineskip
    \begin{subfigure}{0.8\textwidth}
        \includegraphics[width=\textwidth]{../../plots/baseline_vs_llmdynamic_results/llm_weights_evolution_qwen.png}
        \caption{Local Qwen 2.5 7B Meta-Controller (20 interventions)}
    \end{subfigure}
    \caption{Evolution of the Dynamic Reward Weights (post-softmax, as applied to the PBRS signal).}
    \label{fig:llm_weights_evolution}
\end{figure}

\Cref{fig:llm_weights_evolution} reveals strikingly different command strategies. \textbf{Gemini's strategy} as a Commander is characterised by \emph{broad amplification} of the early-chain tasks. $w_{\text{wood}}$ and $w_{\text{stone}}$ rapidly climb and stay high, while $w_{\text{iron}}$ remains near zero after the softmax (mean~0.01). This reveals a counter-intuitive insight: despite the raw LLM outputs assigning $w_{\text{iron}} = 1.0$, the \emph{softmax normalisation} compresses this into a negligible effective weight because the Commander simultaneously assigns high values to wood, stone, and bridge.

\textbf{Qwen's strategy} exhibits a \emph{focused suppression-and-amplification} profile. After an initial exploration phase, $w_{\text{wood}}$ and $w_{\text{stone}}$ are aggressively dropped, while $w_{\text{iron}}$ rises sharply to a peak of 3.87 during the mid-training bottleneck. By zeroing out already-mastered tasks and concentrating the magnetic pull entirely on the bottleneck, Qwen breaks through Iron 30M steps earlier than Gemini.

\subsubsection{Reward Analysis}
\label{sec:exp_rewards}

\begin{figure}[H]
    \centering
    \includegraphics[width=0.7\textwidth]{../../plots/baseline_vs_llmdynamic_results/reward_comparison.png}
    \caption{Average Environment Reward (True Milestones Reached).}
    \label{fig:env_rewards_dynamic}
\end{figure}

\Cref{fig:env_rewards_dynamic} presents the true average environment reward. Both LLM Dynamic runs achieve significantly higher returns than the Baseline, reaching an identical peak reward ($\approx$0.076). This suggests that once the bottleneck is broken, the final quality of the learned policy is comparable; the difference lies purely in the \emph{sample efficiency} driven by Qwen's highly focused interventions.

However, a closer inspection of the training trajectories reveals a critical difference in policy stability and retention. Around step $2.97 \times 10^7$, the Qwen-guided workers exhibit a sharp drop in performance, seemingly forgetting how to complete the Bridge subtask. This catastrophic forgetting is reflected in a corresponding dip in Qwen's average environment reward. Conversely, the Gemini Commander demonstrates superior stability: once it teaches the workers to build a Bridge, they retain this behavior. This retention culminates in a higher average reward spike for Gemini around step $4.65 \times 10^7$. These observations suggest that while Qwen's aggressive, narrowly focused shaping breaks bottlenecks much faster, Gemini's broader amplification strategy (\Cref{fig:llm_weights_evolution}) may foster a more robust, generalizable policy that is less susceptible to catastrophic forgetting.

% ==========================================
% PART 2: HIERARCHICAL REINFORCEMENT LEARNING
% ==========================================
\subsection{Part 2: Hierarchical Reinforcement Learning}
\label{sec:exp_hrl}

While the dynamic shaping architecture successfully guided the lumberjack and miner past initial bottlenecks (like Iron), it ultimately stalled at the combat phase (Enemy, Gold). Modifying the continuous reward landscape was insufficient to teach the workers the discrete, long-horizon macro-actions required for coordinated combat.

This limitation motivated the shift to a full \textbf{Hierarchical Reinforcement Learning (HRL)} architecture, where the Commander does not merely shape the rewards, but explicitly radios discrete \emph{Options} (e.g., ``Craft Sword'', ``Defeat Enemy'') directly to the lower-level MAPPO policy. Three configurations were evaluated over 2000 optimization updates ($\approx$65.5M environment steps each):

\begin{enumerate}
    \item \textbf{No-LoRA Base:} The un-finetuned Qwen 2.5 7B model serving as Meta-Controller.
    \item \textbf{QLoRA (With Reflection):} The QLoRA fine-tuned model with counterfactual reflection logic, triggered when an assigned Option becomes stale (workers fail to complete it within 150 steps).
    \item \textbf{QLoRA (No Reflection):} The identical fine-tuned model, but with the reflection trigger suppressed. On staleness, the Commander re-evaluates the raw environment state without injecting any prior failure context.
\end{enumerate}

\subsubsection{Subtask Progression: Solving the DAG}

\Cref{tab:hrl_breakthrough} presents the exact environment step at which each subtask was first completed across the three configurations. Each one managed to complete the DAG at some point. 

\begin{table}[H]
\centering
\caption{First breakthrough step (millions of environment steps) for each subtask. \textbf{Bold} indicates the fastest configuration per subtask.}
\label{tab:hrl_breakthrough}
\begin{tabular}{lccc}
\hline
\textbf{Subtask} & \textbf{No-LoRA} & \textbf{QLoRA + Refl.} & \textbf{QLoRA -- No Refl.} \\
\hline
Wood       & \textbf{0.16M} & 0.29M & 0.26M \\
Stone      & \textbf{0.16M} & 0.29M & 0.26M \\
Pickaxe    & 13.50M & \textbf{1.25M} & 0.72M \\
Bridge     & 30.61M & \textbf{2.03M} & 1.18M \\
Iron       & 16.61M & \textbf{2.13M} & 3.01M \\
Sword      & 16.71M & \textbf{2.29M} & 3.15M \\
Armor      & 16.78M & \textbf{2.29M} & 3.15M \\
Enemy      & 30.74M & \textbf{2.49M} & 3.24M \\
Gold       & 32.31M & 7.08M & \textbf{6.29M} \\
\hline
\end{tabular}
\end{table}

Several patterns emerge from \Cref{tab:hrl_breakthrough}. First, all three configurations solve the early-game tasks (Wood, Stone) within the first 0.3M steps, with the No-LoRA base marginally faster due to lower inference overhead. Second, the QLoRA models exhibit a dramatic advantage starting at Pickaxe: both fine-tuned runs complete the mid-game milestones (Pickaxe through Armor) approximately \textbf{8--14$\times$ faster} than the base model. The No-LoRA run does not reach Pickaxe until 13.5M steps, nearly $11\times$ slower than QLoRA with Reflection (1.25M), demonstrating that fine-tuning is essential for coherent option sequencing.

Third, the No-LoRA run exhibits a pronounced \emph{ordering anomaly}: it completes Iron, Sword, and Armor ($\sim$16.7M) well before Bridge (30.6M). Since the bridge is a prerequisite for reaching the enemy side of the river, these combat-preparatory milestones accumulate without immediate use. The workers spend $\sim$14M steps fully equipped but unable to cross the river. Both QLoRA runs, by contrast, follow the optimal DAG ordering (building the bridge \emph{before} crafting combat equipment), resulting in a tight cascade from Bridge to Enemy ($<$0.5M steps gap).

Finally, at the Gold milestone, QLoRA without Reflection is the fastest (6.29M), followed closely by QLoRA with Reflection (7.08M). The No-LoRA base reaches Gold at 32.31M, \textbf{5.1$\times$ slower}, confirming that without fine-tuning, the Commander struggles to plan the late-game transition from combat to resource extraction.

\begin{figure}[H]
    \centering
    \begin{subfigure}{0.32\textwidth}
        \includegraphics[width=\textwidth]{../../plots/hrl_reflection_ablation/hrl_ablation_subtask_Wood.png}
        \caption{Wood}
    \end{subfigure}
    \hfill
    \begin{subfigure}{0.32\textwidth}
        \includegraphics[width=\textwidth]{../../plots/hrl_reflection_ablation/hrl_ablation_subtask_Stone.png}
        \caption{Stone}
    \end{subfigure}
    \hfill
    \begin{subfigure}{0.32\textwidth}
        \includegraphics[width=\textwidth]{../../plots/hrl_reflection_ablation/hrl_ablation_subtask_Pickaxe.png}
        \caption{Pickaxe}
    \end{subfigure}
    \vskip 0.5\baselineskip
    \begin{subfigure}{0.32\textwidth}
        \includegraphics[width=\textwidth]{../../plots/hrl_reflection_ablation/hrl_ablation_subtask_Iron.png}
        \caption{Iron}
    \end{subfigure}
    \hfill
    \begin{subfigure}{0.32\textwidth}
        \includegraphics[width=\textwidth]{../../plots/hrl_reflection_ablation/hrl_ablation_subtask_Sword.png}
        \caption{Sword}
    \end{subfigure}
    \hfill
    \begin{subfigure}{0.32\textwidth}
        \includegraphics[width=\textwidth]{../../plots/hrl_reflection_ablation/hrl_ablation_subtask_Armor.png}
        \caption{Armor}
    \end{subfigure}
    \vskip 0.5\baselineskip
    \begin{subfigure}{0.32\textwidth}
        \includegraphics[width=\textwidth]{../../plots/hrl_reflection_ablation/hrl_ablation_subtask_Bridge.png}
        \caption{Bridge}
    \end{subfigure}
    \hfill
    \begin{subfigure}{0.32\textwidth}
        \includegraphics[width=\textwidth]{../../plots/hrl_reflection_ablation/hrl_ablation_subtask_Enemy.png}
        \caption{Enemy}
    \end{subfigure}
    \hfill
    \begin{subfigure}{0.32\textwidth}
        \includegraphics[width=\textwidth]{../../plots/hrl_reflection_ablation/hrl_ablation_subtask_Gold.png}
        \caption{Gold}
    \end{subfigure}
    \caption{Subtask completion rates across all nine environment milestones for the three HRL configurations. The dependency chain flows left-to-right, top-to-bottom.}
    \label{fig:hrl_all_subtasks}
\end{figure}

Crucially, \Cref{fig:hrl_all_subtasks} reveals that \textbf{all three configurations ultimately solve the complete DAG}, reaching 100\% completion on every subtask including Gold. This is a qualitative leap over the Dynamic Shaping architecture (Part~1), which never reached Enemy or Gold. The key differentiator between configurations is not whether they succeed, but \emph{how quickly}: QLoRA with Reflection reaches 100\% on Gold by $\sim$10M steps, QLoRA without Reflection by $\sim$8M steps, while the No-LoRA base requires $\sim$35M steps, consuming more than half the training budget before the workers first collect Gold.

Interestingly, the QLoRA with Reflection run's subtask curves exhibit a distinctive ``staircase'' pattern (\Cref{fig:hrl_all_subtasks}, Bridge and Enemy panels): early plateaus near 0\% are followed by abrupt jumps to 100\%. This indicates that the rapid option cycling caused by the reflection mechanism prevents workers from making incremental progress; instead, the completion rate stays near zero until the PPO policy forces a breakthrough, after which it saturates instantly. By contrast, the QLoRA without Reflection run shows a much smoother, gradual learning ramp (taking roughly 4$\times$ as many updates to transition from 0\% to 90\% on the Enemy subtask compared to No-LoRA). This smooth monotonic ramp is evidence of deliberate, stable option sequencing by the fine-tuned Meta-Controller. The No-LoRA base falls in between, showing delayed but relatively fast transitions once it eventually stumbles upon the correct behavior.

\subsubsection{Reward Trajectories}

\begin{figure}[H]
    \centering
    \includegraphics[width=\textwidth]{../../plots/hrl_reflection_ablation/hrl_ablation_rewards.png}
    \caption{Average Extrinsic Reward across HRL configurations. All three runs converge to comparable final rewards ($\sim$0.326), but the QLoRA runs reach this level $\sim$3$\times$ faster than the No-LoRA base.}
    \label{fig:hrl_rewards}
\end{figure}

The reward trajectories (\Cref{fig:hrl_rewards}) quantify the sample efficiency gap. All three configurations converge to nearly identical final rewards (No-LoRA: 0.326, QLoRA+Refl: 0.326, QLoRA--NoRefl: 0.327), with peak rewards clustered at $\sim$0.335. This convergence confirms that the underlying PPO policy has sufficient capacity to master the environment regardless of the Meta-Controller quality, given enough time. The practical difference is entirely one of sample efficiency.

Compared to the Dynamic Shaping architectures (Part~1), which plateaued at 0.076, the HRL architecture achieves a \textbf{4.4$\times$ improvement} in peak reward. The QLoRA runs reach their peak reward plateau by $\sim$28M steps (No Reflection) and $\sim$37M steps (With Reflection), while the No-LoRA base does not stabilize until $\sim$55M steps. This $\sim$2$\times$ sample efficiency gap between the QLoRA runs is the first quantitative signal that the reflection mechanism, while not preventing convergence, actively delays it.

\subsubsection{Stabilizing Non-Stationary Objectives via KL-Divergence Guarding}
\label{sec:kl_guard}

In the HRL architecture, the transition between discrete macro-actions (Options) induces abrupt shifts in the intrinsic reward landscape. From the perspective of the low-level PPO workers, the objective function is highly non-stationary. To prevent catastrophic forgetting of previously learned kinetic policies, we bound the surrogate objective update using a strict Kullback-Leibler (KL) divergence guard.

During the optimization epochs, we continuously monitor the approximate KL divergence between the behavioral policy $\pi_{\theta_{\text{old}}}$ and the updating policy $\pi_{\theta}$~\parencite{Schulman_Wolski_Dhariwal_Radford_Klimov_2017}:
\begin{equation}
    D_{\text{KL}}^{\text{approx}} = \mathbb{E}\left[ -\log \left( \frac{\pi_{\theta}(a|s, o)}{\pi_{\theta_{\text{old}}}(a|s, o)} \right) \right]
\end{equation}

If $D_{\text{KL}}^{\text{approx}}$ exceeds the predefined trust region threshold ($\delta = 0.015$, implemented as \texttt{target\_kl}), the gradient updates for the current batch are immediately aborted.

\begin{figure}[H]
    \centering
    \includegraphics[width=\textwidth]{../../plots/hrl_reflection_ablation/hrl_ablation_kl_spikes.png}
    \caption{Maximum approximate KL divergence per update across the three HRL configurations. The dashed red line marks the early-stopping threshold $\delta = 0.015$.}
    \label{fig:kl_spikes}
\end{figure}

\begin{table}[H]
\centering
\caption{KL divergence statistics and policy stability metrics across HRL configurations.}
\label{tab:kl_stats}
\begin{tabular}{lccc}
\hline
\textbf{Metric} & \textbf{No-LoRA} & \textbf{QLoRA + Refl.} & \textbf{QLoRA -- No Refl.} \\
\hline
Mean $\max D_{\text{KL}}$       & 0.0066 & 0.0082 & 0.0073 \\
Peak $\max D_{\text{KL}}$       & 0.025  & \textbf{0.117}  & 0.107 \\
KL early-stop triggers           & 46 / 2{,}000 & \textbf{234 / 2{,}000} & 152 / 2{,}000 \\
Staleness resets (total)         & 28{,}927 & \textbf{88{,}711} & 39{,}137 \\
LLM queries dispatched           & 12{,}752 & 12{,}791 & \textbf{13{,}710} \\
Final TD error (last 100 updates)& 0.042  & 0.024  & \textbf{0.023} \\
Mean TD error (all updates)      & 0.046  & 0.060  & 0.048 \\
\hline
\end{tabular}
\end{table}

\Cref{fig:kl_spikes} and \Cref{tab:kl_stats} reveal a counter-intuitive relationship between KL instability and learning dynamics. The QLoRA with Reflection run triggers the KL early-stopping guardrail \textbf{5.1$\times$ more often} than the No-LoRA base (234 vs.\ 46 triggers) and \textbf{1.5$\times$ more often} than the No-Reflection run (234 vs.\ 152). Both QLoRA runs exhibit large peak KL spikes ($>$0.1), indicating aggressive policy shifts during option transitions. The No-LoRA base, by contrast, shows modest peak KL (0.025) and infrequent triggers. Its policy shifts are small because the un-finetuned Commander generates less decisive option switches. 

These aggressive policy shifts and high KL guard triggers in the QLoRA with Reflection run also directly manifest as increased fluctuations in its average extrinsic reward trajectory (\Cref{fig:hrl_rewards}). The constant disruption of the objective function prevents the PPO workers from smoothly optimizing their kinetic behaviors.

Critically, the QLoRA with Reflection run accumulates \textbf{2.3$\times$ more staleness resets} than the No-Reflection run (88{,}711 vs.\ 39{,}137), despite dispatching nearly the same number of LLM queries (12{,}791 vs.\ 13{,}710). This means each reflection-triggered re-query produces an Option that \emph{fails faster}---the Commander actively generates worse plans when told its previous plan failed. The rapid cycling of short-lived Options produces a high frequency of small policy shifts: each individual shift is small enough to avoid the KL guardrail, but the cumulative effect prevents the Critic from converging on a stable value estimate. We term this pattern \emph{death by a thousand cuts}: individually invisible perturbations that collectively devastate sample efficiency.

The No-Reflection run follows the opposite pattern: \emph{infrequent, large, guarded} transitions. The Commander commits to long-horizon Options and rarely switches. When it does (for example, transitioning from \texttt{CRAFT\_SWORD} to \texttt{DEFEAT\_ENEMY}), the policy shift is dramatic (peak KL = 0.107, or $7\times$ the threshold). The guardrail catches these large transitions and aborts the gradient update. Over subsequent updates, the policy absorbs the new objective gradually. This pattern of rare but controlled disruptions produces the lowest mean TD error (0.048) among the three runs.

\subsubsection{The Paradox of Reflection}
\label{sec:reflection_ablation}

The staleness data (\Cref{tab:kl_stats}) provides the mechanistic explanation for the delay caused by reflection. The With-Reflection run dispatches essentially the same number of LLM queries as the other runs ($\sim$12{,}800), yet accumulates \textbf{3.1$\times$ more staleness resets} than the No-LoRA base (88{,}711 vs.\ 28{,}927) and \textbf{2.3$\times$ more} than the No-Reflection run (88{,}711 vs.\ 39{,}137). Each reflection-triggered query thus produces a worse Option, an Option that expires faster and triggers yet another staleness event.

The root cause lies in how the instruction-tuning dataset was constructed. During QLoRA fine-tuning, the dataset synthetically injected counterfactual failures: the training prompt stated \emph{``You previously chose [Incorrect Option]. It failed.''} The model then learned to output the correct Option. Consequently, the Commander associated the reflection trigger with a strong prior: \textbf{my previous choice was logically wrong, and I must switch to a different option.}

During inference, however, option staleness is frequently triggered not by a logical error by the Commander, but by kinetic delays on the ground: the lumberjack navigating around the river terrain, the miner queuing at resource nodes, or simply needing more steps to execute a fundamentally correct plan. When the orchestrator fires the reflection prompt, the Commander, conditioned by its training data, assumes it made a reasoning error. It abandons the correct long-horizon option and switches to an alternative, derailing the workers' execution.

To illustrate this transparency, consider a snapshot directly from the evaluation logs where the workers have already defeated the Enemy and just need to collect the Gold. The Commander correctly assigns \texttt{COLLECT\_GOLD}, but when the workers take too long and trigger a staleness reset, the LLM hallucinates a logical error and regresses to an unnecessary prerequisite:

\begin{verbatim}
Inventory: {..., 'pickaxe': 1, 'sword': 1, 'bridge': 1, 'enemy': 1, 'gold': 0}
Previous Assignment: Agent 1 failed on COLLECT_GOLD

LLM Response:
{
  "reflection": "Previous assignment of IDLE and COLLECT_GOLD failed because 
                 agents lacked the necessary sub-tier items.",
  "dag_check": "A0 waits (bridge done), A1 crafts armor.",
  "agent_0_option": "IDLE",
  "agent_1_option": "CRAFT_ARMOR"
}
\end{verbatim}
Even though Armor is entirely irrelevant once the Enemy is defeated, the Commander doubts its own correct logic because the reflection prompt implied it made a mistake.

We term this failure mode the \emph{Reflection Contradiction}: the architecture inadvertently trained the LLM to doubt its own correct logic. Disabling the reflection mechanism allows the Commander to maintain stable plans. On re-query, it re-confirms its current Option rather than compulsively switching.

\begin{figure}[H]
    \centering
    \includegraphics[width=\textwidth]{../../plots/hrl_reflection_ablation/hrl_ablation_td_error.png}
    \caption{Temporal Difference Error Variance. The With-Reflection configuration exhibits persistent high variance during mid-training from contradictory option switches. Both QLoRA runs converge to similar final TD errors ($\sim$0.024), while the No-LoRA base settles higher at 0.042.}
    \label{fig:hrl_td_error}
\end{figure}

The downstream effect is quantified in the TD Error Variance (\Cref{fig:hrl_td_error}). The With-Reflection configuration maintains the highest mean TD error across the full training run (0.060), 25\% higher than the No-Reflection run (0.048) and 30\% higher than the No-LoRA base (0.046). This elevated TD error during mid-training reflects the Critic's inability to track the rapidly switching intrinsic reward landscape caused by reflection-induced option churn. By the final 100 updates, both QLoRA runs converge to comparable TD errors (With Reflection: 0.024, No Reflection: 0.023), indicating that the Critic eventually stabilizes once the policy has saturated all subtasks and option switching diminishes. The No-LoRA base, despite its lower mid-training variance, settles at a higher final TD error (0.042). Its Critic struggles with the less structured option transitions produced by the un-finetuned model.

Ultimately, all three HRL configurations solve the complete environment DAG, representing a qualitative breakthrough over the Dynamic Shaping architecture. The key finding is that \textbf{QLoRA fine-tuning is essential for sample efficiency} (5$\times$ faster to Gold than the base model), while the reflection mechanism (despite its theoretical appeal) \textbf{delays convergence} by injecting ungrounded self-doubt into the Meta-Controller's reasoning. The KL-divergence guardrail proves indispensable for absorbing the inherent non-stationarity of option transitions, but it cannot protect against the high-frequency, sub-threshold perturbations generated by the reflection loop.

% ==========================================================================
\subsection{Part 3: Robustness and Semantic Adaptation}
\label{sec:robustness_and_adaptation}

Having established that hard integration (HRL) outperforms soft integration (Dynamic Shaping) in solving the complete dependency graph, we now return briefly to the Mixing Board architecture. To fully validate the theoretical claims made in \Cref{sec:ema_smoothing} regarding stability, we isolate two critical variables of the soft integration approach: the numerical stability provided by temporal smoothing, and the zero-shot semantic adaptation capabilities of the LLM Commander.

\subsubsection{Ablation: The Necessity of EMA Smoothing}
\label{sec:ema_ablation}

In \Cref{sec:ema_smoothing}, we posited that feeding raw, discrete priority logits directly into the potential-based shaping function would induce severe non-stationarity. To empirically verify this, an ablation run was conducted wherein the Exponential Moving Average (EMA) layer was disabled.

\begin{figure}[H]
    \centering
    \includegraphics[width=0.48\textwidth]{../../plots/llm_dynamic_ablations_and_shifts/ablation_rewards.png}
    \hfill
    \includegraphics[width=0.48\textwidth]{../../plots/llm_dynamic_ablations_and_shifts/ablation_td_error.png}
    \caption{Comparison of training stability with and without EMA smoothing ($\alpha=0.2$). Disabling smoothing induces a catastrophic spike in Temporal Difference (TD) error variance (right).}
    \label{fig:ema_ablation_td}
\end{figure}

\begin{figure}[H]
    \centering
    \includegraphics[width=0.48\textwidth]{../../plots/llm_dynamic_ablations_and_shifts/ablation_subtask_Iron.png}
    \hfill
    \includegraphics[width=0.48\textwidth]{../../plots/llm_dynamic_ablations_and_shifts/ablation_subtask_Sword.png}
    \caption{Subtask completion rates for Iron (left) and Sword (right). The variance spike directly correlates with suboptimal learning and temporary regressions in advanced subtasks.}
    \label{fig:ema_ablation_subtasks}
\end{figure}

As shown in Figure~\ref{fig:ema_ablation_td}, removing the EMA filter results in a noticeable spike in TD-error variance. The ablation itself was executed by passing a \texttt{skip\_ema=True} flag to the \texttt{LLMOrchestrator}, which bypassed the exponential smoothing step entirely.

The Critic network struggles to track a reward objective that shifts instantaneously. Crucially, Figure~\ref{fig:ema_ablation_subtasks} illustrates the behavioral consequence: the non-stationarity leads to suboptimal optimization, causing the workers to exhibit slight regressions in their \texttt{Iron} and \texttt{Sword} completion rates during the mid-stages of training. However, it is worth noting that the No-EMA workers do eventually recover and climb back to a 100\% success rate on these tasks. Thus, while the workers can mathematically overcome the noise, this proves that the EMA smoothing equation is highly helpful for bridging discrete LLM outputs with continuous RL gradients, enabling faster convergence and preventing such regressions.

To quantify the optimization efficiency provided by the EMA filter, Table~\ref{tab:ema_convergence_speed} details the number of environment steps required to reliably learn (i.e., achieve a $\ge 90\%$ success rate) the advanced subtasks.

\begin{table}[H]
    \centering
    \begin{tabular}{lrr}
        \toprule
        \textbf{Subtask} & \textbf{With-EMA (Baseline)} & \textbf{No-EMA (Ablation)} \\
        \midrule
        \texttt{Wood} & 0.33M steps & 2.36M steps ($7.2\times$ slower) \\
        \texttt{Stone} & 0.82M steps & 2.46M steps ($3.0\times$ slower) \\
        \texttt{Pickaxe} & 1.31M steps & 6.88M steps ($5.3\times$ slower) \\
        \texttt{Bridge} & 2.06M steps & 9.24M steps ($4.5\times$ slower) \\
        \texttt{Iron} & 2.22M steps & 21.00M steps ($9.4\times$ slower) \\
        \texttt{Sword} & 2.52M steps & 21.63M steps ($8.6\times$ slower) \\
        \texttt{Armor} & 2.52M steps & 21.63M steps ($8.6\times$ slower) \\
        \texttt{Enemy} & 3.01M steps & Failed to Converge \\
        \texttt{Gold} & 7.07M steps & Failed to Converge \\
        \bottomrule
    \end{tabular}
    \caption{Environment steps required to reach a 90\% completion rate. Disabling the EMA smoothing filter causes a severe learning slowdown across all tasks, culminating in an inability to converge on the terminal tasks within the training envelope.}
    \label{tab:ema_convergence_speed}
\end{table}

\begin{figure}[H]
    \centering
    \includegraphics[width=0.8\textwidth]{../../plots/llm_dynamic_ablations_and_shifts/ablation_all_subtasks_bar.png}
    \caption{Maximum subtask completion rates achieved across all runs. Despite the early variance spikes, No-EMA eventually masters all tasks identical to the baseline.}
    \label{fig:ema_ablation_all_subtasks}
\end{figure}

\subsubsection{Computational, Financial, and Energy Overheads}
\label{sec:computational_overhead}

A critical practical consideration when integrating LLMs tightly into the RL training loop is the wall-clock time overhead. Because the meta-controller queries the LLM dynamically throughout the episode, relying on external or cloud-based APIs introduces severe network latency and rate-limiting bottlenecks. 

Analyzing the total execution times for a complete 128-update training run reveals this stark reality. The \textbf{Standard HRL Baseline} (which uses a traditional static meta-controller without LLM queries) completed training in \textbf{8.30 hours}. When deploying an external cloud model (Gemini 2.5 Flash-Lite) for the LLM Dynamic run, the training time ballooned to \textbf{12.27 hours}---a roughly 48\% increase purely due to API latency, despite querying a highly optimized "lite" model. Using a heavier state-of-the-art online model (such as GPT-4) was entirely infeasible for our full training comparisons due to strict API rate limits and prohibitive latency over thousands of calls.

In contrast, our optimized \textbf{LLM Dynamic (LoRA)} implementation, running a locally fine-tuned Qwen2.5-7B model, completed the identical training run in just \textbf{7.78 hours}. By keeping inference entirely local, batching requests, and avoiding network overhead, the LLM-guided approach achieved execution speeds completely on par with (and in this instance, marginally faster than) the standard HRL baseline. This proves that dynamic LLM orchestration is not just theoretically advantageous, but computationally viable when deployed efficiently at the edge.

Beyond pure execution speed, the financial and energy footprints of these architectures are critical factors for independent researchers. Relying on commercial online APIs for the meta-controller generates a staggering volume of requests; simulating 128 parallel environments over millions of steps yields millions of input and output tokens per training run. Relying on paid APIs for this volume of queries would render the approach financially prohibitive. Furthermore, reinforcement learning is highly hardware-intensive. A standard local workstation drawing approximately 750W continuously over a 12.27-hour API-bound run---or a 13.5-hour unoptimized local run---results in substantial electricity consumption, which becomes especially burdensome when multiplied by dozens of iterations, failed experiments, and hyperparameter sweeps. In regions with high energy costs such as Switzerland, the ability of the optimized local LoRA implementation to slash training time down to 7.78 hours provides a compounded benefit: it completely eliminates API token costs while simultaneously reducing the electrical footprint and associated utility costs by nearly 40\% per run.

\subsubsection{Environment Shift: The Paradox of Semantic Over-Confidence}
\label{sec:semantic_adaptation}

A fundamental critique of using LLMs in MARL is whether they rely on genuine semantic world-models or merely overfit to node names. To test this, the environment was silently modified at Update 500: the `Workbench` zone was renamed to `Furnace` in the Commander's prompt space. A secondary control run replaced `Workbench` with a nonsense string, `Zylorg`.

This shift was executed by dynamically injecting an alias mapping into the \texttt{LLMOrchestrator} mid-training, forcing the Commander to output weights for the alias rather than the canonical node name, which the orchestrator then reversed and fed to the PPO workers.

The initial hypothesis assumed the Commander would cleanly map "Furnace" to crafting success, while failing on "Zylorg". However, empirical results revealed a counter-intuitive phenomenon: \textbf{semantic over-confidence induced severe reward hacking, while semantic ignorance acted as a natural regularizer.}

\begin{figure}[H]
    \centering
    \includegraphics[width=0.48\textwidth]{../../plots/llm_dynamic_ablations_and_shifts/shift_subtask_Iron.png}
    \hfill
    \includegraphics[width=0.48\textwidth]{../../plots/llm_dynamic_ablations_and_shifts/shift_subtask_Sword.png}
    \caption{Subtask completion rates during semantic shifts. The `Furnace` shift traps workers in a local optimum, halting `Iron` and `Sword` progression, while the `Zylorg` shift organically recovers.}
    \label{fig:semantic_shift_subtasks}
\end{figure}

\begin{figure}[H]
    \centering
    \includegraphics[width=0.8\textwidth]{../../plots/llm_dynamic_ablations_and_shifts/shift_all_subtasks_bar.png}
    \caption{Maximum subtask completion rates for the Shift group. Furnace masters initial tasks (Wood, Stone, Pickaxe) but completely fails advanced tasks requiring exploration (Iron, Sword).}
    \label{fig:shift_all_subtasks}
\end{figure}

As seen in Figures~\ref{fig:semantic_shift_subtasks} and~\ref{fig:shift_all_subtasks}, the `Zylorg` control unexpectedly succeeded across almost all subtasks, achieving near-parity with the baseline's extrinsic reward. When the Commander encountered "Zylorg", it lacked a semantic prior. Consequently, it assigned it a balanced priority weight alongside other nodes, focusing its reasoning on recognizable end-goals like enemies.

\begin{hegquote}
\ttfamily
"reasoning": "Enemy defeat is the current bottleneck with no progress.",\\
"w\_wood": 0.8,\\
"w\_stone": 0.8,\\
"w\_zylorg": 0.9,\\
"w\_iron": 0.9,\\
"w\_enemy": 1.2
\end{hegquote}

Because \texttt{w\_zylorg} was not artificially inflated, the MAPPO workers, driven by the baseline extrinsic rewards for crafting a Sword, naturally explored the environment, discovered the `Zylorg` station was the required prerequisite, and solved the bottleneck via standard PPO exploration. The Commander's ignorance prevented it from interfering with the RL algorithm's foundational mechanics.

Conversely, the `Furnace` shift failed catastrophically on advanced tasks (`Iron`, `Sword`) due to reward hacking. The Commander correctly deduced zero-shot that a Furnace serves an identical semantic purpose to a crafting bench, as demonstrated in its reasoning logs:

\begin{hegquote}
\ttfamily
"reasoning": "Agents are not crafting Pickaxe or Sword, indicating a \\
bottleneck at the Workbench.",\\
"w\_furnace": 0.4,\\
"w\_iron": 0.2
\end{hegquote}

However, by mapping the Furnace to the crafting bottleneck, the Commander heavily skewed the Potential-Based Reward Shaping (PBRS) landscape. The MAPPO workers successfully crafted the Pickaxe (Figure~\ref{fig:shift_all_subtasks}), but subsequently discovered they could accumulate massive intrinsic returns simply by navigating back to the Furnace and idling there. They optimized for the dense shaping signal at the expense of the sparse extrinsic reward, halting all exploration for `Iron` and `Sword`. 

This finding highlights a critical vulnerability in neuro-symbolic MARL architectures: while LLMs possess exceptional zero-shot semantic mapping, coupling high-confidence oracle weights with dense gradient shaping can inadvertently trap low-level workers in local optima.